\scanpage{92}
\mspage{9}
% archive sheet number: 91
\margin{\unsure{\S} On trouve alors une suite exacte de la forme (1.29), mais où on peut écrire $H^i(U_{\bar{\unsure{\eta|y}}},\mathbb Z_\ell)$ au lieu de $H^i(U_{\bar\eta},\mathbb Z_\ell)^P$.}

On trouve alors une suite exacte
\[
\longrightarrow({\unsure{\Phi|\varphi}}^i_\nu)_{\mathrm{tame}}\longrightarrow H^i(X_s,\Lambda_\nu)\longrightarrow H^i(U_{\bar\eta},\Lambda_\nu)^P\longrightarrow(\Phi^{i+1}_\nu)_{\mathrm{tame}}\longrightarrow\cdots
\tag*{(1.29)}
\]
avec
\[
(\Phi^i_\nu)_{\mathrm{tame}}=\varinjlim_{(n,p)=1}H^{\unsure{0|i}}_{Y_n}(X_n,\Lambda_\nu),
\tag*{(1.30)}
\]
où $X_n=X\times_S S(n)$, $S(n)$ le normalisé de $S$ dans $K(n)$, $Y_n$ l'image inverse de $Y=X-U$ dans $X(n)$. La suite spectrale de localisation est ici
\[
(\Phi^*_\nu)_{\mathrm{tame}}\Longleftarrow E_2^{pq}=H^p(X_0,(\underline\Phi^q_\nu)_{\mathrm{tame}})
\tag*{(1.31)}
\]
avec
\[
(\underline\Phi^q_\nu)_{\mathrm{tame}}=\varinjlim_{(n,p)=1}\bigl(\deleted{\unsure{L}}\underline H^q_{Y(n)}((\Lambda_\nu)_{X(n)})\bigr)\mid X(n)_0,
\tag*{(1.32)}
\]
où $X(n)_0$ est la \deleted{le \unsure{revêtement de}} fibre spéciale de $X(n)$, identifiée (comme \deleted{\unsure{objet}}\inserted{site} étale) à $X_0\deleted{\unsure{\simeq}}$.

\inserted{Cette suite \deleted{spectrale}\inserted{exacte} (1.29)\unsure{\S} devient \uline{utile si $P$ opère trivialement} sur $H^i(U_{\bar\eta},\mathbb Z_\ell)$, \deleted{\unsure{néanmoins}} condition qui sera vérifiée en tout cas après extension finie de $S$ (cf \unsure{I}).}
\margin{et celle qui s'en déduit par $\varprojlim_\nu$}
\margin{(\unsure{l'hyp. de déf.|l'hyp. de départ} sur $\mathbb Z_\ell$)}

1.33. Les développements restent valables \inserted{mutatis mutandis} si on remplace $P$ par n'importe quel sous-groupe invariant fermé $Q$ de $I$ \deleted{tel que $I$} qui est ``premier à $\ell$'', i.e. tel que $I/Q$ couvre le quotient $\mathbb Z_\ell(1)$ de $I_t$, composante $\ell$-primaire de $I/P$. Un cas particulier \deleted{\unsure{intéressant}} qui peut être utile est celui où $I/Q\simeq\mathbb Z_\ell(1)$ est cette composante $\ell$-primaire, ce qui signifie qu'on fait parcourir à $K_\alpha$ les \deleted{\unsure{de}} $K_{\unsure{\alpha}}(n)$ avec $n$ de la forme $\ell^r$. \deleted{Dans ce cas}
